\section{Dependency Graph}

\begin{figure}[H]
\centering
\begin{tikzpicture}[
  >=Latex,
  box/.style={
    draw,
    rounded corners,
    align=left,
    text width=3.4cm,
    font=\small,
    inner sep=4pt
  },
  imp/.style={box, fill=blue!6},
  core/.style={box, fill=green!6},
  res/.style={box, fill=orange!8},
  ex/.style={box, fill=purple!7}
]
\node[imp]  (tools)  at (0, 0)   {Imported tools\\Wasserstein gradient flows\\Propagation of chaos\\Analytic zero-set facts};
\node[core] (setup)  at (0,-3.2) {Setup\\two-layer model, empirical measure $\hat\rho^{(N)}$, mean-field risk $R(\rho)$, DD / noisy DD};
\node[core] (assump) at (-6.0,-6.4) {Assumptions\\A1--A3 for PDE limit\\A4 for strong noisy-DD analysis};
\node[res]  (prop1)  at (-2.0,-6.4) {Proposition 1\\finite-width optimum is within $O(1/N)$ of $\inf_\rho R(\rho)$};
\node[res]  (thm3)   at (2.0,-6.4) {Theorem 3\\SGD / noisy SGD $\to$ PDE / diffusion PDE on finite time windows};
\node[res]  (prop23) at (6.0,-6.4) {Propositions 2--3\\characterize fixed points\\and Boltzmann fixed point of noisy DD};
\node[res]  (thm4)   at (2.0,-9.8) {Theorem 4\\unique minimizer of free energy\\and convergence of diffusion PDE};
\node[res]  (thm67)  at (6.0,-9.8) {Theorems 6--7\\local stability of atomic fixed points\\and exclusion of unstable mixed fixed points};
\node[res]  (thm5)   at (2.0,-13.2) {Theorem 5\\regularized noisy SGD reaches near-optimal risk};
\node[ex]   (lemma1) at (-4.0,-16.6) {Lemma 1\\criterion for spherical-uniform minimizer in isotropic model};
\node[ex]   (thm12)  at (2.0,-16.6) {Theorems 1--2\\application to isotropic / anisotropic Gaussian examples};

\draw[->] (tools) -- (setup);
\draw[->] (setup) -- (prop1);
\draw[->] (setup) -- (thm3);
\draw[->] (tools) -- (thm3);
\draw[->] (assump) -- (thm3);
\draw[->] (assump) -- (prop23);
\draw[->] (thm3) -- (thm4);
\draw[->] (prop23) -- (thm4);
\draw[->] (thm4) -- (thm5);
\draw[->] (prop1) -- (thm5);
\draw[->] (prop23) -- (thm67);
\draw[->] (setup) -- (lemma1);
\draw[->] (lemma1) -- (thm12);
\draw[->] (thm3) -- (thm12);
\draw[->] (prop1) -- (thm12);
\end{tikzpicture}
\caption{Main proof spine of \cite{mei2018meanfield}. The central transition is from finite-width SGD to distributional dynamics, after which the analysis moves to fixed points, free energy, and example-specific PDE reductions. A few logically true but visually redundant edges are omitted so the main proof spine remains readable.}
\end{figure}

\paragraph{How to read the graph.}
The paper's most important move is Theorem~3: once SGD is approximated by DD, the optimization problem is replaced by the study of a measure-valued flow. The noisy branch (Proposition~3 $\to$ Theorem~4 $\to$ Theorem~5) gives the cleanest generic convergence result. The noiseless branch (Propositions~2, Theorems~6--7) gives a weaker but structurally informative picture of stable and unstable fixed points.
